\section{How this report was produced}\label{app:method}

This report was written with heavy AI assistance. A reader should know how,
because the way a document is made affects how far it can be trusted.

\paragraph{Division of work.} One coordinating model (Claude Opus) planned
the report, fixed its structure and terminology, and settled conflicts between
workers. Separate worker agents did the evidence work, each in its own
workstream: a PoC evidence audit, scientific foundations, adjacent work and
novelty, interpretation of the N2 and N3 results, the implemented
architecture, the future platform, the validation program, the organizational
path, the education path, the program plan, and the figures. Each worker
wrote notes to \repo{report/notes/} before any section was drafted. The notes
are kept with the report so that every section can be traced to them.

\paragraph{Evidence rules.} Every PoC number was checked in the primary
artifact (a run ledger, a gap-map JSON file, code, or a committed report),
not copied from a summary. Where a committed summary and a raw artifact
disagreed, the raw artifact won, and the disagreement is stated in the text.
The test counts in this report come from running the four offline test
suites during the audit. Numbers that could not be traced to a committed
data file (for example the 3-of-19 precision figure) are reported with that
qualification.

\paragraph{Citation rules.} Every reference in the bibliography was checked
against Crossref, OpenAlex, arXiv or the publisher page (title, authors, year,
venue, DOI or stable URL) by the agent that proposed it. The repository's own
research notes were treated as leads, not proof. Several leads failed this
check and were dropped or corrected; for example, the widely repeated claim
that experts omit about 70\% of what they know has no single primary study
behind it, so the report cites the primary omission studies instead
(\cref{sec:problem}).

\paragraph{Reviews.} The complete draft went through two review rounds, each
with two independent reviewers who did not see each other's findings:
\begin{itemize}[nosep]
  \item Round 1: a scientific reviewer (research logic, literature
    positioning, unsupported claims, novelty) and a technical reviewer
    (correspondence with the code, architecture accuracy, accidental claims
    that future components exist).
  \item Round 2: a hostile grant reviewer (why it matters, what is new,
    whether the hypothesis is falsifiable and the validation credible) and a
    clarity and reproducibility reviewer.
\end{itemize}
Both rounds are complete, and the report was revised after each one. Round 1 changed the
absence-check wording, the claims about span provenance, the gap map's display cap, the gate
logic, and cut duplicated material. Round 2 restructured the validation program around a
pre-grant V1 closure study, tied each money tranche to a gate reading CONTINUE, powered V5 to
decide H2, moved V3 (E-OSS) to the organizational track as a separate hypothesis that can neither
kill nor license H2, added an external gate reader, added citations a novelty check had missed,
and ran this clarity pass. A final audit then checked major statements against literature,
repository evidence, or an explicit hypothesis or future-work label. The figures were redrawn
after a visual inspection found several printing below 7~pt from over-wide scaling, and the
compiled PDF was inspected page by page.

\paragraph{Limits of this process.} The workers and the reviewers are
models. Agreement between them is not evidence. The checks above reduce, but
do not remove, the risk of errors in reading sources. The decisions on scope,
direction and the research program remain the project owner's.
